%% Paper 2 — Springer Nature Proceedings (sn-jnl). BAN TIENG VIET (de doc/duyet).
%%
%% >>> Bien dich bang pdfLaTeX. Tieng Viet co dau nho babel + [T5]{fontenc}.
%%     Neu dau bi loi -> doi Compiler sang XeLaTeX va thay fontenc/babel bang:
%%       \usepackage{fontspec}
%%       \usepackage{polyglossia}\setmainlanguage{vietnamese}
\documentclass[pdflatex,sn-mathphys-num]{sn-jnl}

\usepackage{lmodern}
\usepackage[utf8]{inputenc}
\usepackage[T5]{fontenc}
\usepackage[vietnamese]{babel}

\addto\captionsvietnamese{%
  \renewcommand{\abstractname}{Abstract}%
  \renewcommand{\tablename}{Table}%
  \renewcommand{\figurename}{Fig.}%
  \renewcommand{\refname}{References}%
  \renewcommand{\bibname}{References}%
}

\usepackage{graphicx}
\graphicspath{{figures/}{./}{../figures/}}
\usepackage{multirow}
\usepackage{amsmath,amssymb,amsfonts}
\usepackage{booktabs}

\raggedbottom

\begin{document}

\title[Bộ đếm nhân nhẹ count-only kèm bất định hiệu chuẩn]{Bộ đếm nhân tế bào nhẹ, chỉ giám sát bằng
số đếm, kèm định lượng bất định được hiệu chuẩn cho ảnh mô bệnh học}

\author[1]{\fnm{Thi Thu Hiep} \sur{Dinh}}\email{dinhthithuhiep1@gmail.com}
\author*[1]{\fnm{Viet Hang} \sur{Duong}}\email{hangdv@uit.edu.vn}

\affil*[1]{\orgdiv{Khoa Khoa học Máy tính}, \orgname{Trường Đại học Công nghệ Thông tin (UIT), ĐHQG-HCM},
\orgaddress{\city{TP. Hồ Chí Minh}, \country{Việt Nam}}}

\abstract{%
Đếm nhân tế bào trên ảnh mô bệnh học nhuộm H\&E là bước định lượng nền tảng cho nhiều tác vụ chẩn đoán, nhưng
phần lớn phương pháp hiện nay đòi hỏi mặt nạ phân vùng ở mức điểm ảnh --- loại nhãn đắt và tốn công chuyên gia.
Chúng tôi đặt câu hỏi: để đếm, liệu có thật sự cần mask không? Bài báo trình bày một bộ đếm nhân nhẹ (3,62
triệu tham số) chỉ huấn luyện bằng một con số đếm cho mỗi ảnh (count-only), không dùng mask, kèm đầu định
lượng bất định được hiệu chuẩn (neo-Poisson $\sigma$ và conformal prediction). Qua thí nghiệm có kiểm soát (cố
định kiến trúc, chỉ thay đổi loại giám sát), chúng tôi cho thấy: (i) giám sát count-only đạt độ chính xác đếm
ngang bằng giám sát bằng mask trên cùng kiến trúc, trong khi nhãn rẻ hơn một bậc; (ii) mô hình nhẹ này vượt ba
dòng phương pháp nhãn-đắt (density, instance segmentation, detection) ngay trên bài đếm; (iii) chỉ 25\% nhãn đã
đạt khoảng 93\% hiệu năng đầy đủ; (iv) đầu bất định cho khoảng tin cậy hiệu chuẩn cả ở mức biên lẫn theo từng
cơ quan, và chuyển giao coverage sang dữ liệu ngoài miền. Kết quả nhất quán trên ba tập dữ liệu độc lập
(PanNuke, NuInsSeg, MoNuSeg).
}

\keywords{đếm nhân tế bào, học nhãn-hiệu quả, giám sát yếu, định lượng bất định, conformal prediction,
mô bệnh học tính toán}

\maketitle

%%===========================================================%%
\section{Giới thiệu}\label{sec:intro}

Mật độ và số lượng nhân tế bào là chỉ dấu định lượng quan trọng trong bệnh học: chúng liên quan tới mức độ tăng
sinh, phân độ khối u và đánh giá vi môi trường mô. Các phương pháp học sâu hiện đại chủ yếu đi theo hướng phân
vùng thực thể (instance segmentation) hoặc hồi quy bản đồ mật độ (density map), và đều cần nhãn mask ở mức điểm
ảnh --- mỗi nhân phải được khoanh viền thủ công. Chi phí gán nhãn này là rào cản lớn khi mở rộng sang mô hoặc
cơ quan mới chưa có chú thích.

Trong khi đó, ở bài toán đếm đám đông (crowd counting) người ta đã dùng giám sát chỉ-đếm (count-only) --- chỉ
cần tổng số đối tượng cho mỗi ảnh. Tuy nhiên trong đếm nhân mô bệnh học, chuẩn mực vẫn là mask hoặc point. Câu
hỏi trung tâm của bài báo: liệu giám sát count-only có đủ để đếm nhân chính xác không, và ta được hay mất gì so
với giám sát bằng mask?

Chúng tôi tiếp cận theo tinh thần thực nghiệm có kiểm soát: cố định kiến trúc, chỉ thay đổi loại giám sát, để
tách bạch đóng góp của giám sát khỏi đóng góp của kiến trúc. Bài báo không tuyên bố một kiến trúc mới (bộ mã
hóa EfficientNet-lite0 và bộ giải mã U-Net là tổ hợp tiêu chuẩn); đóng góp nằm ở bằng chứng thực nghiệm và gói
ứng dụng:

\begin{enumerate}
\item Chứng minh count-only xấp xỉ mask trên cùng kiến trúc (Mục~\ref{sec:parity}): độ chính xác đếm tương
đương, nhưng nhãn rẻ hơn một bậc.
\item So sánh bốn paradigm công bằng, cùng tập kiểm tra leak-free (Mục~\ref{sec:main}): bộ đếm nhẹ count-only
vượt cả density, instance và detection; đặc biệt detection sụp đổ ở vùng đông đúc --- bằng chứng trực tiếp cho
động lực crowding.
\item Định lượng hiệu quả-nhãn (Mục~\ref{sec:label}): đường cong học theo tỉ lệ nhãn.
\item Đầu bất định hiệu chuẩn nhẹ (Mục~\ref{sec:uq}): tham số hóa $\sigma$ neo-Poisson và conformal cho
coverage đúng mục tiêu ở mức biên và theo cơ quan, chuyển giao được sang dữ liệu ngoài miền.
\end{enumerate}

Toàn bộ số liệu được báo cáo trên một mô hình duy nhất (EfficientNet-lite0, 3,62M tham số) với giao thức
leak-free.

%%===========================================================%%
\section{Công trình liên quan}\label{sec:related}

\textbf{Phân vùng và đếm nhân.} Các phương pháp SOTA (HoVer-Net, StarDist, CellViT và biến thể trên PanNuke)
đạt độ chính xác cao nhưng dựa vào mask điểm ảnh; đếm thường là readout từ phân vùng. Gần đây, một số công
trình chưng cất tri thức từ foundation model bệnh học xuống mạng nhỏ để đếm (ví dụ chưng cất từ H-Optimus
xuống khoảng 24M tham số, báo cáo $R^2$ đếm khoảng 0,871) nhưng không kèm định lượng bất định.

\textbf{Giám sát count-only.} Phổ biến ở crowd counting, hiếm ở đếm nhân mô bệnh học. Chúng tôi khảo sát trực
tiếp giá trị của loại giám sát này trong miền bệnh học.

\textbf{Định lượng bất định.} Conformal prediction cho khoảng tin cậy phi-phân phối với bảo đảm coverage. Với
dữ liệu đếm (phương sai tăng theo kỳ vọng), tham số hóa $\sigma$ hợp lý là then chốt; chúng tôi dùng dạng
neo-Poisson $\sigma = \sqrt{\max(\mu,1)}\cdot\exp(\log s)$.

\textbf{Kiến trúc.} Bộ mã hóa EfficientNet-lite và bộ giải mã U-Net là mẫu chuẩn. Chúng tôi không coi kiến trúc
là đóng góp; nó là nền được kiểm soát để so sánh giám sát.

%%===========================================================%%
\section{Phương pháp}\label{sec:method}

\subsection{Kiến trúc}
Mô hình \texttt{DensitySigmaUNet}: bộ mã hóa EfficientNet-lite0 (tiền huấn luyện ImageNet) và bộ giải mã kiểu
U-Net, hai đầu ra. Đầu mật độ suy ra số đếm $\mu = \sum \text{density}$ (tổng toàn ảnh). Đầu bất định
$\sigma = \sqrt{\max(\mu,1)}\cdot\exp(\mathrm{clamp}(\log s,-2,2))$ theo dạng neo-Poisson, học phần tán xạ trên
nền phương sai Poisson; $\mu$ được tách gradient khỏi nhánh $\sigma$. Cấu hình $ch{=}32$ cho 3,62 triệu tham số.

\subsection{Giám sát chỉ-đếm}
Hàm mất mát: $\mathcal{L} = w_d\mathcal{L}_{density} + w_c\mathcal{L}_{count} + w_{nll}\mathcal{L}_{nll}$.
Count-only đặt $w_d{=}0, w_c{=}1$ (chỉ học từ số đếm mức ảnh $|\mu - \mathrm{GT}|$, không dùng mask). Baseline
mask đối chứng đặt $w_d{=}1, w_c{=}0{,}01$ (bản đồ mật độ mục tiêu dựng từ mask GT). $\mathcal{L}_{nll}$ là
$\beta$-NLL ($\beta{=}0,5$) huấn luyện $\sigma$; $\mu$ được detach.

\subsection{Hiệu chuẩn bất định bằng conformal}
Sau huấn luyện, dùng split-conformal trên tập hiệu chuẩn để suy khoảng $[\mu - q\sigma, \mu + q\sigma]$ đạt mức
phủ mục tiêu $1-\alpha$ ($\alpha{=}0,1$). Ba sơ đồ: global (một ngưỡng $q$), mondrian ($q$ riêng theo cơ quan),
cluster ($q$ theo nhóm độ khó).

%%===========================================================%%
\section{Thiết lập thực nghiệm}\label{sec:setup}

\textbf{Tập dữ liệu (độc lập, đa nguồn).} PanNuke (pan-cancer, $256^2$, 19 loại mô) --- in-domain chính, train
fold~1+2 sang test fold~3 (leak-free); NuInsSeg (đa cơ quan H\&E) --- in-domain thứ hai; MoNuSeg~2018 (đa cơ
quan TCGA, $1000^2$ sang patch~256 native) --- in-domain thứ ba, độc lập nguồn; CryoNuSeg (cryosection) --- phân
tích giới hạn.

\textbf{Giao thức leak-free.} Với PanNuke, huấn luyện fold~1+2, kiểm tra fold~3; mọi mô hình so sánh đều huấn
luyện lại từ đầu trên cùng tập, chấm trên cùng tập kiểm tra. Mỗi cấu hình chạy ít nhất 3 seed.

\textbf{Chỉ số.} Đếm: $R^2$, MAE. Bất định: coverage biên, coverage cơ quan tệ nhất (worst-organ), độ rộng
khoảng (width), điểm Winkler.

%%===========================================================%%
\section{Kết quả}\label{sec:results}

\subsection{So sánh bốn paradigm (PanNuke, leak-free)}\label{sec:main}

\begin{table}[t]
\centering
\caption{So sánh bốn paradigm trên PanNuke (leak-free, cùng tập kiểm tra).}\label{tab:main}
\begin{tabular}{@{}lccccc@{}}
\toprule
Mô hình & Tham số & Nhãn & Paradigm & $R^2\uparrow$ & MAE$\downarrow$ \\
\midrule
\textbf{Ours --- count-only} & \textbf{3,62M} & đếm & hồi quy & \textbf{0,920} & \textbf{3,61} \\
StarDist & 1,43M & mask & instance & 0,841 & 4,64 \\
DualU-Net & $\sim$35M & mask & density & 0,804 & 5,33 \\
DeNuC (2026) & 4,37M & box/point & detection & 0,455 & 9,64 \\
\botrule
\end{tabular}
\end{table}

Bộ đếm nhẹ count-only vượt cả ba baseline nhãn-đắt trên chính bài đếm, với nhãn rẻ hơn một bậc. Đáng chú ý,
detection sụp đổ đơn điệu theo mật độ (độ lệch từ $-0,8$ ở vùng thưa xuống $-30,4$ ở vùng $\ge 45$ nhân) --- khi
nhân chồng lấn, hậu xử lý NMS bóp chết dương-tính-thật nên recall sụp. Đây là bằng chứng trực tiếp cho động lực
crowding. Về tham số, mô hình của chúng tôi nhỏ hơn DualU-Net khoảng 10 lần nhưng lớn hơn StarDist; luận điểm
là chính xác nhất, nhãn rẻ nhất và có định lượng bất định, không phải nhỏ nhất.

\subsection{Count-only so với mask, trên cùng kiến trúc}\label{sec:parity}

\begin{table}[t]
\centering
\caption{Cố định kiến trúc efflite0, chỉ đổi loại giám sát (PanNuke, leak-free).}\label{tab:parity}
\begin{tabular}{@{}lcc@{}}
\toprule
Giám sát & $R^2\uparrow$ & MAE$\downarrow$ \\
\midrule
\textbf{count-only} & \textbf{0,9296 $\pm$ 0,009} & \textbf{3,48 $\pm$ 0,19} \\
GT-density (mask) & 0,9247 $\pm$ 0,009 & 3,58 $\pm$ 0,16 \\
$\Delta$ & $+0,005$ & $-0,10$ \\
\botrule
\end{tabular}
\end{table}

Chênh lệch ($+0,005$ $R^2$) nhỏ hơn độ lệch chuẩn ($\pm 0,009$), nên về mặt thống kê là ngang nhau (parity).
Kết luận trung thực: không phải count-only tốt hơn mask, mà là count-only đạt độ chính xác tương đương mask
trên cùng kiến trúc, trong khi chỉ cần một con số đếm cho mỗi ảnh thay vì khoanh viền pixel. Đây là đối chứng
đóng chốt cho định vị bài báo.

\subsection{Hiệu quả-nhãn}\label{sec:label}

\begin{table}[t]
\centering
\caption{Đường cong học theo tỉ lệ nhãn (PanNuke, 3 seed).}\label{tab:label}
\begin{tabular}{@{}lccccc@{}}
\toprule
\% nhãn & 5\% & 10\% & 25\% & 50\% & 100\% \\
\midrule
$R^2$ & 0,45$\pm$0,21 & 0,74 & \textbf{0,85$\pm$0,03} & 0,86 & 0,92$\pm$0,01 \\
\botrule
\end{tabular}
\end{table}

Chỉ 25\% nhãn đã đạt khoảng 93\% hiệu năng đầy đủ; 100\% khớp headline 0,920 (nhất quán).

\subsection{Hiệu quả tính toán}
Mô hình 3,62M tham số, 0,98 GMACs ($256^2$), độ trễ 6,0\,ms (GPU) và 46\,ms (CPU), peak VRAM 37,9\,MB ---
thời-gian-thực, chạy được cả trên GPU khiêm tốn lẫn CPU.

\subsection{Tổng quát hóa đa tập và định lượng bất định}\label{sec:uq}

\begin{table}[t]
\centering
\caption{In-domain trên ba tập độc lập (cùng giao thức count-only).}\label{tab:gen}
\begin{tabular}{@{}lcc@{}}
\toprule
Tập & $R^2$ & MAE \\
\midrule
PanNuke & 0,920 $\pm$ 0,005 & 3,61 \\
NuInsSeg & 0,854 $\pm$ 0,048 & --- \\
MoNuSeg (official split) & 0,663 $\pm$ 0,099 & 5,42 $\pm$ 0,86 \\
\botrule
\end{tabular}
\end{table}

\begin{table}[t]
\centering
\caption{Conformal in-domain (PanNuke, 20 seed).}\label{tab:conf}
\begin{tabular}{@{}lccccc@{}}
\toprule
Sơ đồ & cov biên & width & Winkler & worst-org & \#thiếu phủ \\
\midrule
global & 0,904 & 14,5 & 20,3 & 0,765 & 4/19 \\
\textbf{mondrian} & 0,923 & 16,6 & 20,5 & \textbf{0,906} & \textbf{0/19} \\
cluster & 0,906 & 14,8 & \textbf{20,0} & 0,859 & 0/19 \\
\botrule
\end{tabular}
\end{table}

Coverage biên đạt mục tiêu 0,90; mondrian nâng coverage cơ quan tệ nhất lên 0,906 (0/19 cơ quan thiếu phủ).
Thêm đầu $\sigma$ tốn 0 độ chính xác (nhờ detach $\mu$).

\begin{table}[t]
\centering
\caption{Ablation $\sigma$ neo-Poisson so với $\sigma$ thô (cùng mô hình, chỉ đổi đầu $\sigma$).}\label{tab:sigma}
\begin{tabular}{@{}llcccc@{}}
\toprule
Tập (dải count) & $\sigma$ & cov & width & Winkler & worst-org \\
\midrule
PanNuke (0--45, hẹp) & Poisson & 0,905 & 14,1 & 24,7 & 0,64 \\
                     & thô     & 0,908 & 15,1 & \textbf{21,6} & 0,66 \\
\textbf{NuInsSeg (0--500+, rộng)} & \textbf{Poisson} & 0,926 & \textbf{37,7} & \textbf{45,4} & \textbf{0,86} \\
                     & thô     & 0,937 & 67,9 & 78,1 & 0,66 \\
\botrule
\end{tabular}
\end{table}

Trên dải count hẹp, hai dạng $\sigma$ tương đương (thậm chí $\sigma$ thô nhỉnh). Trên dải count rộng, neo-Poisson
cắt 44\% width, 42\% Winkler và nâng worst-organ $+0,21$ cùng mức phủ, vì $\sigma$ hằng số buộc phải phồng
khoảng cho mọi ảnh. Kết luận: neo-Poisson cần thiết khi dải count rộng (tình huống triển khai thực tế), thừa
khi hẹp.

%%===========================================================%%
\section{Thảo luận và Hạn chế}\label{sec:discuss}

\textbf{Chuyển giao ngoài miền.} Huấn luyện một tập rồi kiểm tra zero-shot trên tập khác: coverage của khoảng
tin cậy chuyển giao được (sau hiệu chuẩn lại ở đích), nhưng độ chính xác điểm (MAE) không chuyển giao khi dải
count hoặc độ phóng đại lệch nhiều. Trường hợp cực đoan CryoNuSeg (count tới 761 vượt tầm huấn luyện) cho thấy
giới hạn: $\mu$ under-predict mạnh, coverage chỉ giữ được bằng khoảng rất rộng. Điểm tích cực: bảo đảm coverage
đến từ tầng conformal (phi-phân phối), không phải từ $\sigma$; độ rộng khoảng chính là tín hiệu đã ra ngoài
vùng tin cậy.

\textbf{Phụ thuộc scale nhân.} Đầu mật độ nhạy với kích thước nhân theo điểm ảnh; khi thu nhỏ ảnh quá mạnh,
$\mu$ có thể sụp. Ta xử lý bằng cách cắt patch ở scale native.

\textbf{Về đóng góp.} Kiến trúc là tổ hợp có sẵn --- bài báo không tuyên bố kiến trúc mới. Giá trị nằm ở thiết
kế thí nghiệm có kiểm soát, bằng chứng count-only đủ để đếm, và đầu bất định hiệu chuẩn mà các bộ đếm nén hiện
có chưa có.

%%===========================================================%%
\section{Kết luận}\label{sec:conclusion}

Chúng tôi cho thấy một bộ đếm nhân nhẹ (3,62M), chỉ giám sát bằng số đếm mức ảnh, đạt độ chính xác đếm ngang
giám sát bằng mask trên cùng kiến trúc và vượt các baseline nhãn-đắt ngay trên bài đếm, với nhãn rẻ hơn một
bậc. Mô hình kèm đầu bất định hiệu chuẩn (neo-Poisson và conformal) cho coverage đúng mục tiêu ở mức biên và
theo cơ quan, và chuyển giao coverage sang dữ liệu ngoài miền. Hướng phát triển: khắc phục phụ thuộc scale bằng
học đa độ phân giải, và mở rộng đầu bất định để bắt cả sai số do dịch chuyển miền.

%%===========================================================%%
\backmatter

\bmhead{Declarations}
Đây là công trình gốc, chưa được công bố và không đồng thời được nộp ở nơi khác. Các tác giả tuyên bố không có
xung đột lợi ích.

%%===========================================================%%
\bibliography{references}

\end{document}
